\chapter{Bayesian Methods}\label{ch:qm:bayesian-methods}

A multi-manager platform reviews its fifty portfolio managers. The best three-year \omterm{def:qm:estimation:sharpe}{Sharpe ratio} is 2.4, and the allocation committee
proposes to double that manager's capital. Three years estimate a \omterm{def:qm:estimation:sharpe}{Sharpe ratio} with a \omterm{def:qm:estimation:estimator}{standard error} of about 0.58, and the best of fifty
noisy records is, almost by construction, a good manager who was also lucky. Shrunk toward what the platform's managers usually achieve,
the same record is worth a \omterm{def:qm:estimation:sharpe}{Sharpe ratio} of 1.14; the manager's next three years, in the chapter's simulation, deliver 1.69, and across two
thousand simulated platforms the best record averages 2.09 while the same manager's next three years average 1.02. This chapter is about
reasoning with prior knowledge: priors, posteriors and the conjugate families that make them easy; shrinkage and the \omterm{def:qm:estimation:estimator}{estimator} that
showed the sample mean can be beaten; \omterm{def:qm:bayesian-methods:hier}{hierarchical models} whose prior is learnt from the population; and the Markov-chain samplers that
compute posteriors when no formula does.

\section{Priors, posteriors and conjugacy}

\begin{definition}[Prior, posterior, conjugate prior]\label{def:qm:bayesian-methods:prior}
In a Bayesian model the parameter $\theta$ is random. Its \emph{prior distribution}\index{prior distribution} $p(\theta)$ states what is
believed before the data; the \emph{posterior distribution}\index{posterior distribution} $p(\theta \mid x) \propto p(x \mid \theta)\,
p(\theta)$ is the Bayesian update of it by the likelihood of the data $x$ (Bayes, 1763). A family of priors is a \emph{conjugate
prior}\index{conjugate prior} for a likelihood if the posterior stays in the family.
\end{definition}

\begin{definition}[Posterior predictive distribution, credible interval]\label{def:qm:bayesian-methods:predictive}
The \emph{posterior predictive distribution}\index{posterior predictive distribution} of new data $\tilde x$ is $p(\tilde x \mid x) = \int
p(\tilde x \mid \theta)\,p(\theta \mid x)\,d\theta$. A $(1 - \alpha)$ \emph{credible interval}\index{credible interval} is an interval to
which the posterior gives probability $1 - \alpha$, usually between its $\alpha/2$ and $1 - \alpha/2$ quantiles.
\end{definition}

Unlike a confidence interval, a \omterm{def:qm:bayesian-methods:predictive}{credible interval} is a probability statement about the parameter given the data, and it is only as good as
the prior. Three conjugate pairs cover most desk uses. A hit rate $p$ with a $\mathrm{Beta}(a, b)$ prior, after $k$ successes in $n$ trials,
has a $\mathrm{Beta}(a + k, b + n - k)$ posterior: a signal right on 110 of 200 trades, with a $\mathrm{Beta}(50, 50)$ prior that says
signals are rarely far from a coin, has posterior mean 0.533, a 95\% \omterm{def:qm:bayesian-methods:predictive}{credible interval} $(0.477, 0.589)$ and a posterior probability of
0.876 that its hit rate exceeds one half. A normal mean with a normal prior gives a normal posterior (below). A normal sample with
unknown mean and variance has the normal-inverse-gamma family.

\begin{proposition}[The normal-normal posterior is a precision-weighted average]\label{prop:qm:bayesian-methods:normal}
If $\theta \sim \mathcal N(m, \tau^2)$ and $x \mid \theta \sim \mathcal N(\theta, \sigma^2)$, then $\theta \mid x \sim \mathcal N(\hat\theta,
v)$ with $1/v = 1/\tau^2 + 1/\sigma^2$ and
\[
\hat\theta = \frac{m/\tau^2 + x/\sigma^2}{1/\tau^2 + 1/\sigma^2} = m + (1 - B)(x - m), \qquad B = \frac{\sigma^2}{\sigma^2 + \tau^2}.
\]
\end{proposition}

\begin{proof}
The log posterior is $-\frac{(\theta - m)^2}{2\tau^2} - \frac{(x - \theta)^2}{2\sigma^2}$ up to a constant, a quadratic in $\theta$ with
coefficient $-\frac12(1/\tau^2 + 1/\sigma^2)$ on $\theta^2$ and $m/\tau^2 + x/\sigma^2$ on $\theta$; completing the square gives the mean
and variance.
\end{proof}

The shrinkage factor $B$ is the share of the observed deviation $x - m$ that is attributed to noise. It is large when the data are noisy
($\sigma$ large) or the population homogeneous ($\tau$ small).

\section{Shrinkage}

\begin{definition}[Shrinkage estimator, James--Stein estimator]\label{def:qm:bayesian-methods:js}
A \emph{shrinkage estimator}\index{shrinkage estimator} pulls a set of estimates toward a common target. For $X \sim \mathcal N(\theta,
\sigma^2I_d)$, the \emph{James--Stein estimator}\index{James--Stein estimator} toward zero is $\bigl(1 - (d - 2)\sigma^2/\lvert X\rvert^2\bigr)X$;
toward the grand mean $\bar X$ it is $\bar X + \bigl(1 - (d - 3)\sigma^2/\sum_i(X_i - \bar X)^2\bigr)(X - \bar X)$, and its positive part
truncates the factor at zero.
\end{definition}

\begin{theorem}[James--Stein dominates the sample mean]\label{thm:qm:bayesian-methods:js}
For $d \ge 3$ and every $\theta$, the \omterm{def:qm:bayesian-methods:js}{James--Stein estimator} toward zero has a strictly smaller total mean squared error than $X$:
$\E\lvert\delta_{\mathrm{JS}} - \theta\rvert^2 = d\sigma^2 - (d - 2)^2\sigma^4\,\E[1/\lvert X\rvert^2] < d\sigma^2$.
\end{theorem}

\begin{proof}
Take $\sigma = 1$ and write $\delta = X - g(X)$ with $g(x) = (d - 2)x/\lvert x\rvert^2$. Then $\E\lvert\delta - \theta\rvert^2 = d -
2\E[(X - \theta)\cdot g(X)] + \E\lvert g(X)\rvert^2$. Stein's identity, integration by parts against the normal density, gives
$\E[(X_i - \theta_i)g_i(X)] = \E[\partial_ig_i(X)]$, and $\sum_i\partial_ig_i(x) = (d - 2)^2/\lvert x\rvert^2$. Since $\lvert g(x)\rvert^2 = (d
- 2)^2/\lvert x\rvert^2$ too, the risk is $d - (d - 2)^2\E[1/\lvert X\rvert^2]$, and the expectation is finite for $d \ge 3$.
\end{proof}

The result (James and Stein, 1961) shocked a profession: estimating three unrelated means, the sample means are inadmissible. The gain
is large when the means are close to the target relative to the noise, which is the situation of fifty managers' \omterm{def:qm:estimation:sharpe}{Sharpe ratios} measured
over three years. Efron and Morris (1975) made it famous with batting averages.

\section{Hierarchical models and empirical Bayes}

\begin{definition}[Hierarchical model, empirical Bayes]\label{def:qm:bayesian-methods:hier}
A \emph{hierarchical model}\index{hierarchical model} gives the parameters of many units a common prior whose own parameters are
unknown: $x_i \mid \theta_i \sim \mathcal N(\theta_i, \sigma_i^2)$, $\theta_i \sim \mathcal N(m, \tau^2)$. \emph{Empirical
Bayes}\index{empirical Bayes} estimates $(m, \tau^2)$ from the data, by maximising the marginal likelihood $x_i \sim \mathcal N(m, \sigma_i^2 +
\tau^2)$ or by moments, and then applies the posterior formulas of each unit with the estimates plugged in.
\end{definition}

With equal $\sigma_i = \sigma$, the moment estimates are $\hat m = \bar x$ and $\hat\tau^2 = \max(0, s_x^2 - \sigma^2)$: the spread of the
records beyond what noise alone would produce. The resulting \omterm{def:qm:estimation:estimator}{estimator} is James--Stein's up to the factor $(k - 3)/(k - 1)$, which explains
the theorem: shrinkage toward the grand mean is Bayes with a prior estimated from the other units.

\begin{example}[Fifty managers]\label{ex:qm:bayesian-methods:platform}
The platform's fifty managers have true annual \omterm{def:qm:estimation:sharpe}{Sharpe ratios} drawn from $\mathcal N(0.5, 0.4^2)$, and each three-year record adds noise of
standard deviation $\sigma = 1/\sqrt3 = 0.577$ (chapter~11). On the seeded platform, whose best record is 2.41 as in the hook, the records
have mean 0.51 and standard deviation 0.71, so $\hat\tau = \sqrt{0.71^2 - 1/3} = 0.41$ and $B = 0.333/(0.333 + 0.168) = 0.665$: two thirds of
every deviation from the mean is noise. The best record becomes $0.51 + 0.335 \times (2.41 - 0.51) = 1.14$, with posterior standard deviation
0.33; James--Stein gives 1.20. The second and third records, 1.82 and 1.45, become 0.95 and 0.83; the ranking does not change, only the
distances. The true ratios of the top three are 1.07, 0.96 and 1.68, and their next three years deliver 1.69, 0.57 and 0.99
(\cref{fig:qm:bayesian-methods:shrink}).
\end{example}

\begin{omfigure}
\begin{tikzpicture}
\begin{axis}[width=9.5cm,height=6.2cm,xlabel={Sharpe ratio, first three years},ylabel={Sharpe ratio, next three years},
  tick label style={font=\small},label style={font=\small},xmin=-1,xmax=3,ymin=-1,ymax=3,grid=major,grid style={gray!25},
  legend columns=3,legend style={font=\footnotesize,at={(0.5,-0.25)},anchor=north,draw=none,fill=none,
  /tikz/every even column/.append style={column sep=8pt}},table/col sep=comma]
\addplot[only marks,mark=*,mark size=1.6pt,omDef] table[x=x1,y=x2]{figdata/methods/14-bayesian-methods/managers.csv};
\addplot[black,dashed,thick] table[x=x,y=diag]{figdata/methods/14-bayesian-methods/lines.csv};
\addplot[omThm,very thick] table[x=x,y=eb]{figdata/methods/14-bayesian-methods/lines.csv};
\legend{managers,no shrinkage,empirical Bayes}
\end{axis}
\end{tikzpicture}

\omcaption{The fifty managers' \omterm{def:qm:estimation:sharpe}{Sharpe ratios} over two consecutive three-year periods. Taking the first record at face value predicts the
dashed diagonal; the empirical-Bayes posterior mean predicts the flatter line, slope 0.335 through the platform mean 0.51, and the second
period scatters around it. Data: the chapter's tutorial, seeded.}\label{fig:qm:bayesian-methods:shrink}
\end{omfigure}

Measured against the true ratios, \omterm{def:qm:bayesian-methods:hier}{empirical Bayes} cuts the mean squared error of the fifty estimates from 0.335 to 0.149 on this platform.
Averaged over 2\,000 simulated platforms it cuts it from 0.336 to 0.121, close to the 0.109 of the oracle that knows $m$ and $\tau$; and
against the next three years' records, which is what an allocator can check, from 0.672 to 0.457, a reduction of 32\%
(\cref{fig:qm:bayesian-methods:mse}).

\begin{omfigure}
\begin{tikzpicture}
\begin{axis}[width=10cm,height=5cm,ybar,bar width=12pt,
  xtick={0,1,2,3},xticklabels={{raw\\record},{James--\\Stein},{empirical\\Bayes},{oracle\\Bayes}},
  xticklabel style={align=center},
  ylabel={mean squared error},tick label style={font=\small},label style={font=\small},
  xmin=-0.6,xmax=3.6,ymin=0,ymax=0.8,grid=major,grid style={gray!25},
  legend style={font=\footnotesize,at={(0.97,0.97)},anchor=north east,fill=white,draw=none},table/col sep=comma]
\addplot[fill=omDef!60,draw=omDef] table[x=k,y=truth]{figdata/methods/14-bayesian-methods/mse.csv};
\addplot[fill=omThm!45,draw=omThm] table[x=k,y=next]{figdata/methods/14-bayesian-methods/mse.csv};
\legend{against the true ratios,against the next three years}
\end{axis}
\end{tikzpicture}

\omcaption{Mean squared error of four estimates of fifty managers' \omterm{def:qm:estimation:sharpe}{Sharpe ratios}, averaged over 2\,000 simulated platforms: the raw
three-year records (0.336 against the truth, 0.672 against the next period), James--Stein and \omterm{def:qm:bayesian-methods:hier}{empirical Bayes} (0.121 and 0.457), and the
posterior mean with the true platform parameters (0.109 and 0.445). Data: the chapter's tutorial, seeded.}\label{fig:qm:bayesian-methods:mse}
\end{omfigure}

The best record is where shrinkage matters most. Across the 2\,000 platforms, the best of fifty three-year records averages 2.09 while
the same manager's true \omterm{def:qm:estimation:sharpe}{Sharpe ratio} averages 1.01: the selection alone adds 1.08. The empirical-Bayes estimate of the selected manager
averages 1.01, and the next three years 1.02 (\cref{fig:qm:bayesian-methods:curse}). The committee's problem is a winner's curse, the
selection effect of chapter~12's search, and the prior is its correction.

\begin{omfigure}
\begin{tikzpicture}
\begin{axis}[width=10cm,height=5cm,xlabel={estimate minus the true Sharpe ratio, best of fifty},ylabel={density},
  tick label style={font=\small},label style={font=\small},xmin=-1.5,xmax=3,ymin=0,ymax=1.6,grid=major,grid style={gray!25},
  legend style={font=\footnotesize,at={(0.97,0.97)},anchor=north east,fill=white,draw=none},table/col sep=comma]
\addplot[omThm,thick,const plot mark mid,fill=omThm!15] table[x=x,y=raw]{figdata/methods/14-bayesian-methods/curse.csv};
\addplot[omDef,very thick,const plot mark mid] table[x=x,y=eb]{figdata/methods/14-bayesian-methods/curse.csv};
\addplot[black,dashed,thin,sharp plot] coordinates {(0,0) (0,1.6)};
\legend{raw record,empirical Bayes}
\end{axis}
\end{tikzpicture}

\omcaption{The error in the estimated \omterm{def:qm:estimation:sharpe}{Sharpe ratio} of the best of fifty managers, over 2\,000 simulated platforms: the raw record
overstates the chosen manager by 1.08 on average; the empirical-Bayes estimate is centred on the truth. Data: the chapter's tutorial,
seeded.}\label{fig:qm:bayesian-methods:curse}
\end{omfigure}

The \omterm{def:qm:bayesian-methods:predictive}{posterior predictive distribution} of the best manager's next three-year \omterm{def:qm:estimation:sharpe}{Sharpe ratio} is normal with mean 1.14 and standard deviation
$\sqrt{0.33^2 + 1/3} = 0.67$: a 4.3\% chance of a negative period, and a 2.9\% chance of matching the 2.41 that got him noticed.

\section{Markov chain Monte Carlo}

\omterm{def:qm:bayesian-methods:hier}{Empirical Bayes} plugs in $\hat\tau$ as if it were known; with fifty units $\tau$ is uncertain, and a full Bayesian treatment integrates
over it. The posterior is then no longer normal, and is computed by simulation.

\begin{definition}[Markov chain Monte Carlo, Metropolis--Hastings, Gibbs sampler]\label{def:qm:bayesian-methods:mcmc}
\emph{Markov chain Monte Carlo}\index{Markov chain Monte Carlo} (MCMC) draws from a density $\pi$ known up to a constant by running a
\omterm{def:qm:markov-chains-and-queues:ctmc}{Markov chain} whose \omterm{def:qm:stochastic-differential-equations:kolmogorov}{stationary distribution} is $\pi$. The \emph{Metropolis--Hastings algorithm}\index{Metropolis--Hastings algorithm}
proposes $y \sim q(\cdot \mid x)$ from the current state $x$ and accepts it with probability $\alpha(x, y) = \min\bigl(1,
\frac{\pi(y)q(x \mid y)}{\pi(x)q(y \mid x)}\bigr)$, staying at $x$ otherwise; with a symmetric random-walk proposal the ratio is
$\pi(y)/\pi(x)$. The \emph{Gibbs sampler}\index{Gibbs sampler} updates one block of coordinates at a time by a draw from its conditional
law given the others.
\end{definition}

\begin{proposition}[Detailed balance of Metropolis--Hastings]\label{prop:qm:bayesian-methods:mh}
The Metropolis--Hastings kernel satisfies $\pi(x)K(x, y) = \pi(y)K(y, x)$ for $x \ne y$, so $\pi$ is stationary. A Gibbs update is a
Metropolis--Hastings step whose proposal is the conditional law, and its acceptance probability is 1.
\end{proposition}

\begin{proof}
For $x \ne y$, $\pi(x)K(x, y) = \pi(x)q(y \mid x)\alpha(x, y) = \min\bigl(\pi(x)q(y \mid x), \pi(y)q(x \mid y)\bigr)$, which is symmetric in
$(x, y)$. Summing \omterm{def:qm:markov-chains-and-queues:balance}{detailed balance} over $x$ gives $\sum_x\pi(x)K(x, y) = \pi(y)$ (the discrete-time analogue of chapter~8's
condition). For Gibbs, if $x$ and $y$ differ only in block $j$, $q(y \mid x) = \pi(y_j \mid x_{-j})$ and $x_{-j} = y_{-j}$, so
$\pi(y)q(x \mid y) = \pi(x_{-j})\pi(y_j \mid x_{-j})\pi(x_j \mid x_{-j}) = \pi(x)q(y \mid x)$, and the ratio is 1.
\end{proof}

\begin{definition}[Burn-in, effective sample size]\label{def:qm:bayesian-methods:ess}
The \emph{burn-in}\index{burn-in} is the initial stretch of a chain discarded because it still depends on the starting point. The
\emph{effective sample size}\index{effective sample size} of $n$ correlated draws is $n/(1 + 2\sum_{k \ge 1}\rho_k)$, $\rho_k$ the chain's
autocorrelations: the number of independent draws with the same variance of the mean.
\end{definition}

For the \omterm{def:qm:bayesian-methods:hier}{hierarchical model} with flat priors on $m$ and $\tau$, every conditional is standard: each $\theta_i$ is normal (the proposition
above), $m$ is normal around $\bar\theta$, and $\tau^2$ is inverse gamma. A \omterm{def:qm:bayesian-methods:mcmc}{Gibbs sampler} of 20\,000 sweeps (2\,000 of \omterm{def:qm:bayesian-methods:ess}{burn-in}) gives the
best manager a posterior mean of 1.12 with a 95\% \omterm{def:qm:bayesian-methods:predictive}{credible interval} $(0.41, 2.01)$, wider than \omterm{def:qm:bayesian-methods:hier}{empirical Bayes}' $1.14 \pm 1.96 \times 0.33$
because $\tau$ is uncertain: its posterior mean is 0.40 with a 95\% interval $(0.13, 0.66)$. A random-walk Metropolis sampler on the
marginal posterior of $(m, \ln\tau)$, four chains from dispersed starts, accepts 50\% of its proposals and agrees (posterior mean of $\tau$
0.39, \cref{fig:qm:bayesian-methods:tau}); its Gelman--Rubin statistic is 1.001, and the 18\,000 kept draws of $\tau$ are worth 506
independent ones from the Gibbs chain and 767 from one Metropolis chain.

\begin{omfigure}
\begin{tikzpicture}
\begin{axis}[width=10cm,height=5cm,xlabel={$\tau$, dispersion of true Sharpe ratios},ylabel={posterior density},
  tick label style={font=\small},label style={font=\small},xmin=0,xmax=1,ymin=0,ymax=3.6,grid=major,grid style={gray!25},
  legend columns=3,legend style={font=\footnotesize,at={(0.5,-0.3)},anchor=north,draw=none,fill=none,
  /tikz/every even column/.append style={column sep=8pt}},table/col sep=comma]
\addplot[omThm,thick,const plot mark mid,fill=omThm!15] table[x=tau,y=gibbs]{figdata/methods/14-bayesian-methods/tau.csv};
\addplot[omDef,very thick,const plot mark mid] table[x=tau,y=metropolis]{figdata/methods/14-bayesian-methods/tau.csv};
\addplot[black,dashed,thin,sharp plot] coordinates {(0.41,0) (0.41,3.6)};
\legend{Gibbs sampler,Metropolis (four chains),empirical Bayes $\hat\tau$}
\end{axis}
\end{tikzpicture}

\omcaption{Posterior of the dispersion $\tau$ of the platform's true \omterm{def:qm:estimation:sharpe}{Sharpe ratios} from the \omterm{def:qm:bayesian-methods:mcmc}{Gibbs sampler} (18\,000 draws) and from four
random-walk Metropolis chains on the marginal posterior, against the empirical-Bayes point estimate 0.41 (dashed). Data: the chapter's
tutorial, seeded.}\label{fig:qm:bayesian-methods:tau}
\end{omfigure}

\section{Tutorial: the allocator's shortlist}\label{tut:qm:bayesian-methods:shortlist}

\begin{tutorial}
\textbf{Goal.} Shrink fifty managers' \omterm{def:qm:estimation:sharpe}{Sharpe ratios} by \omterm{def:qm:bayesian-methods:hier}{empirical Bayes} and by a \omterm{def:qm:bayesian-methods:mcmc}{Gibbs sampler}, and check the shrinkage on the next three
years. \textbf{End state:} \cref{fig:qm:bayesian-methods:shrink,fig:qm:bayesian-methods:mse,fig:qm:bayesian-methods:tau}, the best
manager's 2.41 shrunk to 1.14, and the 32\% cut in out-of-sample error.
\begin{tutsteps}
\item \textbf{\omterm{def:qm:bayesian-methods:hier}{Empirical Bayes}} for normal means with equal or unequal \omterm{def:qm:estimation:estimator}{standard errors}.
\omcode{code/firm/bayes/firm_bayes.py}{45}{66}{Empirical-Bayes shrinkage of normal means.}\label{lst:qm:bayesian-methods:eb}
\item \textbf{The \omterm{def:qm:bayesian-methods:mcmc}{Gibbs sampler}} for the hierarchical normal model.
\omcode{code/firm/bayes/firm_bayes.py}{96}{113}{Gibbs sampler for $x_i \sim \mathcal N(\theta_i, \sigma_i^2)$, $\theta_i \sim \mathcal N(m,
\tau^2)$.}\label{lst:qm:bayesian-methods:gibbs}
\item \textbf{Run} \texttt{problem()}, \texttt{average\_mse()}, \texttt{mcmc()} in \texttt{qm\_bayes.py}, then \texttt{fig\_bayes.py}.
\end{tutsteps}
\textbf{What to change next.} Give the managers different track-record lengths (unequal $\sigma_i$) and watch the short records shrink
more; replace the normal prior by a Student $t$ so that a genuinely exceptional manager is shrunk less, sampling it with the Metropolis
routine.
\end{tutorial}

\section{Build: the Bayesian toolkit}\label{bld:qm:bayesian-methods:bayes}

\begin{build}
\bfield{Purpose} Every ranking of many noisy estimates in the miniature firm (managers, signals, venues, counterparties) is shrunk before it
is acted on; this module does the shrinking and the sampling.
\bfield{Interface} \texttt{beta\_binomial}, \texttt{normal\_normal}, \texttt{nig\_update}; \texttt{eb\_normal\_means(x, se)} returning the
prior's $m$ and $\tau^2$, the shrinkage factors and the posterior means and standard deviations; \texttt{james\_stein(x, se)};
\texttt{metropolis(logpdf, x0, n, step, rng)}; \texttt{gibbs\_hierarchical(x, se, n\_iter, rng)}; \texttt{ess(chain)};
\texttt{rhat(chains)}.
\bfield{Rules} Priors and their parameters are reported with every posterior; samplers take a generator and return the whole chain, and no
posterior summary is reported without an \omterm{def:qm:bayesian-methods:ess}{effective sample size} and, for several chains, $\hat R$.
\bfield{Acceptance tests} \texttt{code/firm/bayes/tests/}: the conjugate formulas; \omterm{def:qm:bayesian-methods:hier}{empirical Bayes} recovers $(m, \tau)$ and its fixed point
maximises the marginal likelihood with unequal errors; James--Stein beats the raw estimates; Metropolis samples a normal; the \omterm{def:qm:bayesian-methods:ess}{effective
sample size} of an AR(1) chain is $n(1 - \phi)/(1 + \phi)$; Gibbs agrees with \omterm{def:qm:bayesian-methods:hier}{empirical Bayes} when the units are many.
\bfield{Stretch} A Student-$t$ prior by Metropolis-within-Gibbs; Hamiltonian Monte Carlo on the marginal posterior; the Robbins formula for
Poisson counts.
\end{build}

\begin{omsources}
\item T.~Bayes, ``An essay towards solving a problem in the doctrine of chances'', \emph{Philosophical Transactions of the Royal Society}
53, 1763.
\item W.~James and C.~Stein, ``Estimation with quadratic loss'', \emph{Fourth Berkeley Symposium}, vol.~1, 1961.
\item B.~Efron and C.~Morris, ``Data analysis using Stein's estimator and its generalizations'', \emph{Journal of the American Statistical
Association} 70, 1975.
\item N.~Metropolis, A.~W.~Rosenbluth, M.~N.~Rosenbluth, A.~H.~Teller and E.~Teller, ``Equation of state calculations by fast computing
machines'', \emph{Journal of Chemical Physics} 21, 1953.
\item W.~K.~Hastings, ``Monte Carlo sampling methods using Markov chains and their applications'', \emph{Biometrika} 57, 1970.
\item S.~Geman and D.~Geman, ``Stochastic relaxation, Gibbs distributions, and the Bayesian restoration of images'', \emph{IEEE
Transactions on Pattern Analysis and Machine Intelligence} 6, 1984.
\item A.~Gelman and D.~B.~Rubin, ``Inference from iterative simulation using multiple sequences'', \emph{Statistical Science} 7, 1992.
\end{omsources}

\section{Exercises}

\begin{exercise}[$\star$]\label{exo:qm:bayesian-methods:1}
A new signal was right on 30 of 50 trades. With a $\mathrm{Beta}(50, 50)$ prior, what are the posterior and its mean? With a flat
$\mathrm{Beta}(1, 1)$ prior?
\end{exercise}

\begin{exercise}[$\star$]\label{exo:qm:bayesian-methods:2}
A manager's five-year \omterm{def:qm:estimation:sharpe}{Sharpe ratio} is 1.5. With a prior $\mathcal N(0.5, 0.4^2)$ for true \omterm{def:qm:estimation:sharpe}{Sharpe ratios} and a \omterm{def:qm:estimation:estimator}{standard error} of
$1/\sqrt5$, what are the posterior mean and standard deviation?
\end{exercise}

\begin{exercise}[$\star$]\label{exo:qm:bayesian-methods:3}
A chain of 10\,000 draws has lag-one autocorrelation 0.8 and behaves like an AR(1). What is its \omterm{def:qm:bayesian-methods:ess}{effective sample size}?
\end{exercise}

\begin{exercise}[$\star\star$]\label{exo:qm:bayesian-methods:4}
Show that with $k$ units, equal $\sigma$ and the moment estimate $\hat\tau^2 = s_x^2 - \sigma^2 > 0$, the empirical-Bayes shrinkage factor is
$\sigma^2/s_x^2$, and compare it with the James--Stein factor toward the grand mean.
\end{exercise}

\begin{exercise}[$\star\star$]\label{exo:qm:bayesian-methods:5}
In the platform example, how many years of track record would make the shrinkage factor $B$ fall to one half?
\end{exercise}

\begin{exercise}[$\star\star$]\label{exo:qm:bayesian-methods:6}
Derive the conditional law of $\tau^2$ in the \omterm{def:qm:bayesian-methods:mcmc}{Gibbs sampler} when the prior on $\tau$ is flat.
\end{exercise}

\begin{exercise}[$\star\star\star$]\label{exo:qm:bayesian-methods:7}
\emph{Coding.} Over 2\,000 simulated platforms, compare the average best three-year record, the chosen manager's true \omterm{def:qm:estimation:sharpe}{Sharpe ratio}, its
empirical-Bayes estimate and its next three years.
\end{exercise}

\begin{exercise}[$\star\star\star$]\label{exo:qm:bayesian-methods:8}
\emph{Find the flaw.} ``We shrink every signal's backtest \omterm{def:qm:estimation:sharpe}{Sharpe ratio} toward zero by the James--Stein factor, so our estimates are
conservative and we no longer need to worry about how many signals we tried.''
\end{exercise}

\section{Problem: The Allocator's Shortlist}

\begin{problem}[{Weekend problem --- shrinking the best of fifty managers}]\label{pb:qm:bayesian-methods:1}
A platform has fifty managers whose true annual \omterm{def:qm:estimation:sharpe}{Sharpe ratios} are drawn from $\mathcal N(0.5, 0.4^2)$; each has a three-year record with
\omterm{def:qm:estimation:estimator}{standard error} $1/\sqrt3$ and, later, a second three-year record. The platform is the chapter's seeded one.

\medskip\noindent\textbf{Part I --- The records.}
\begin{enumerate}
\item What are the best record, and the mean and standard deviation of the fifty records?
\item How many managers have a negative record?
\item What standard deviation would the records have if every manager had the same true \omterm{def:qm:estimation:sharpe}{Sharpe ratio}?
\item What are the moment estimates of $m$ and $\tau$?
\item What is the shrinkage factor $B$?
\end{enumerate}

\medskip\noindent\textbf{Part II --- Shrinking.}
\begin{enumerate}[resume]
\item What are the best manager's posterior mean and standard deviation?
\item What does James--Stein give?
\item What do the second and third records become, and does the ranking change?
\item What were the top three's true \omterm{def:qm:estimation:sharpe}{Sharpe ratios}, and their next three years?
\item What is the posterior predictive probability that the best manager's next three years are negative?
\end{enumerate}

\medskip\noindent\textbf{Part III --- Checking it.}
\begin{enumerate}[resume]
\item By how much does \omterm{def:qm:bayesian-methods:hier}{empirical Bayes} cut the error against the true ratios on this platform?
\item And on average over 2\,000 platforms, against the truth and against the next three years?
\item On average, what are the best record, the chosen manager's true ratio and its shrunk estimate?
\item What do the \omterm{def:qm:bayesian-methods:mcmc}{Gibbs sampler}'s posterior mean and 95\% interval for the best manager say?
\item What is the posterior of $\tau$, and why does it matter?
\end{enumerate}

\medskip\noindent\textbf{Part IV --- Judgement.}
\begin{enumerate}[resume]
\item Should the committee double the best manager's capital?
\item What prior would you use for a manager hired from outside the platform?
\item When would the normal prior shrink a genuinely exceptional manager too much?
\item State the \emph{named result}: the best manager's posterior mean and the out-of-sample reduction in error.
\item In one sentence: why does the best of many noisy records need shrinking?
\end{enumerate}
\end{problem}

\section{Interview questions}

\begin{interviewq}[$\star$ \iqroles{researcher, trader}]\label{iq:qm:bayesian-methods:1}
What is the difference between a confidence interval and a \omterm{def:qm:bayesian-methods:predictive}{credible interval}?
\end{interviewq}

\begin{interviewq}[$\star\star$ \iqroles{researcher, trader}]\label{iq:qm:bayesian-methods:2}
You must pick one of fifty managers by their three-year \omterm{def:qm:estimation:sharpe}{Sharpe ratios}. How do you estimate what the best one will deliver?
\end{interviewq}

\begin{interviewq}[$\star\star$ \iqroles{researcher, mle}]\label{iq:qm:bayesian-methods:3}
Why can the \omterm{def:qm:bayesian-methods:js}{James--Stein estimator} beat the sample means, and what is the catch?
\end{interviewq}

\begin{interviewq}[$\star\star$ \iqroles{researcher, mle}]\label{iq:qm:bayesian-methods:4}
How do you know an MCMC sampler has converged?
\end{interviewq}

\begin{interviewq}[$\star\star$ \iqroles{mle}]\label{iq:qm:bayesian-methods:5}
How does ridge regression relate to a Bayesian prior?
\end{interviewq}

\begin{interviewq}[$\star\star\star$ \iqroles{researcher}]\label{iq:qm:bayesian-methods:6}
Why does Metropolis--Hastings work, and why is a Gibbs step always accepted?
\end{interviewq}
